\chapter{Core Framework and Graphics API Reference}
\label{ch:api-quick-reference}

This appendix is an orientation map to the public core and graphics surface at the pinned CNA
revision.  It intentionally groups overload families instead of reproducing declarations: the
module-owned headers remain the signature authority, while the cited chapters own behavioral and
renderer-specific qualifications.  A listed method proves API presence only; it does not imply
that every one of the \RendererFamilyCount{} implementation families renders it with the same semantics.

\section{Math and Geometry Types}
\label{sec:api-math}

\begin{small}
\begin{longtable}{>{\raggedright\arraybackslash}p{3cm}>{\raggedright\arraybackslash}p{5.25cm}>{\raggedright\arraybackslash}p{3.25cm}}
\toprule
\textbf{Type group} & \textbf{Public shape} & \textbf{Boundary to remember} \\
\midrule
\endfirsthead
\toprule
\textbf{Type group} & \textbf{Public shape} & \textbf{Boundary to remember} \\
\midrule
\endhead
\cnaclass{Vector2}, \cnaclass{Vector3}, \cnaclass{Vector4} & Public components, constants,
  arithmetic, distance/dot, interpolation, normalization, reflection, and matrix/quaternion
  transforms; \cnaclass{Vector3} adds cross product. & Value-returning and output-reference
  overloads are independent API rows and need independent tests. \\
\cnaclass{Matrix} & Sixteen public elements, direction/translation properties, decomposition,
  arithmetic, inversion/transposition, and the XNA creation families. & CNA uses row-vector
  conventions; \texttt{ToColumnMajor} is a marked bridge for shader upload, not a change of the
  public convention. \\
\cnaclass{Quaternion} & Construction from axis/matrix/yaw-pitch-roll, arithmetic,
  concatenate, conjugate, inverse, lerp and slerp. & Quaternion multiplication order and
  normalization are behavioral contracts, not naming details. \\
\cnaclass{Point}, \cnaclass{Rectangle} & Integer screen-space coordinates, containment,
  intersection/union, offset and inflation. & Rectangle edges and empty/intersection cases use
  exact XNA/FNA boundary rules. \\
\cnaclass{Color} & Packed RGBA value, named colours, vector/float/integer constructors,
  interpolation and premultiplication helpers. & Packed storage is \texttt{AABBGGRR}; constructor
  type changes can alter overload and clamping behavior. \\
\cnaclass{Plane}, \cnaclass{Ray}, \cnaclass{BoundingBox}, \cnaclass{BoundingSphere},
\cnaclass{BoundingFrustum} & Shape construction, containment, intersection, corners and
  transform-related queries. & C\texttt{++} optionals/output references translate nullable and
  \texttt{out} forms; compare semantics, not signature strings alone. \\
\bottomrule
\end{longtable}
\end{small}

Chapter~\ref{ch:math-core-types} covers value semantics and Chapter~\ref{ch:math-conventions}
covers coordinate, matrix, viewport and vertex-layout conventions.  The xna4-spec comparison in
Chapter~\ref{ch:xna4-spec-auditing} shows why generated XML is useful inventory but requires
language normalization and behavioral oracles.

\section{Game Framework and Content}
\label{sec:api-game-content}

\begin{small}
\begin{longtable}{>{\raggedright\arraybackslash}p{2.8cm}>{\raggedright\arraybackslash}p{5.45cm}>{\raggedright\arraybackslash}p{3.25cm}}
\toprule
\textbf{Type} & \textbf{Primary surface} & \textbf{Current qualification} \\
\midrule
\endfirsthead
\toprule
\textbf{Type} & \textbf{Primary surface} & \textbf{Current qualification} \\
\midrule
\endhead
\cnaclass{Game} & Lifecycle events; components, content, services, window and device; fixed/variable
  timing; \texttt{Run}, \texttt{Tick}, \texttt{Exit}, and the initialize/load/update/draw hooks. &
  Host ownership differs on native, Emscripten and mobile paths; the call order is the contract,
  not merely the presence of virtual methods. \\
\cnaclass{GameTime} & Total and elapsed time plus the slow-frame flag. & Mutable by the framework;
  user code consumes the snapshot. \\
\cnaclass{GameWindow} & Bounds, title, resize/orientation/display events and screen-device-change
  protocol, plus marked SDL/native-window conveniences. & CNA uses one SDL3-backed concrete
  class; lifecycle/platform behavior remains host-specific. \\
\cnaclass{GraphicsDeviceManager} & Preferred backbuffer/depth/profile/fullscreen settings,
  \texttt{ApplyChanges}, device creation/preparation events, and presentation-mode extensions. &
  Stored presentation values are requested state, not proof the driver/window accepted them. \\
\cnaclass{ContentManager} & Root directory, service provider, generic \texttt{Load<T>},
  \texttt{Unload}, reader/CNJ registration, manifest and reader-usage extensions. & Resolution is
  XNB first, then a registered loose reader that may select CNJ or a native format; cache and
  ownership vary by asset type. \\
\bottomrule
\end{longtable}
\end{small}

The exact game-loop order is in Chapter~\ref{ch:game-loop}.  Chapters~\ref{ch:content-and-assets}
through~\ref{ch:cnj-format} separate path resolution, XNB container/object graph, and CNJ; a
successful lookup in one route is not evidence for the others.

\section{Graphics Core: Resources and Batching}
\label{sec:api-graphics-core}

\begin{small}
\begin{longtable}{>{\raggedright\arraybackslash}p{3cm}>{\raggedright\arraybackslash}p{5.15cm}>{\raggedright\arraybackslash}p{3.35cm}}
\toprule
\textbf{Type group} & \textbf{Public shape} & \textbf{Evidence boundary} \\
\midrule
\endfirsthead
\toprule
\textbf{Type group} & \textbf{Public shape} & \textbf{Evidence boundary} \\
\midrule
\endhead
\cnaclass{Texture2D} & Dimensions/bounds, whole and rectangular \texttt{SetData}/\texttt{GetData},
  stream decode, PNG/JPEG save, and marked pixel helpers. & A round trip can observe CNA's CPU
  shadow rather than sampled GPU memory; mip/subregion behavior is renderer-specific. \\
\cnaclass{Texture3D}, \cnaclass{TextureCube} & Volume or face dimensions and level/subregion
  upload/readback families. & Construction, format, mip, upload and readback support are separate
  capabilities; no CPU shadow provides a universal fallback. \\
\cnaclass{RenderTarget2D}, \cnaclass{RenderTargetCube} & Texture surface plus depth/stencil,
  multisample and usage metadata; cube adds face selection. & Bind, preserve/discard, plural
  targets, cube unbind and readback have independent renderer paths. \\
\cnaclass{VertexBuffer}, \cnaclass{IndexBuffer} & Counts/usage/declaration plus typed and sliced
  upload/readback; dynamic subclasses add \texttt{SetDataOptions}. & The constructor's dynamic
  intent, upload hint, CPU shadow and renderer buffer are four distinct layers. \\
\cnaclass{SpriteBatch} & \texttt{Begin}/\texttt{End}, texture \texttt{Draw} overloads and
  \texttt{DrawString} for strings/StringBuilder. & Sort mode, state objects, custom effect,
  transform and clipping must reach the selected renderer; a visible sprite alone does not prove
  every argument. \\
\cnaclass{SpriteFont} & Atlas, glyph/cropping/character/kerning arrays, spacing/default character,
  and measurement. & XNB, CNJ and hand-built fonts enter through different reader/ownership
  routes. \\
\bottomrule
\end{longtable}
\end{small}

Chapters~\ref{ch:textures-rendertargets} and~\ref{ch:spritebatch} own the detailed behavior.
Appendix~\ref{ch:feature-matrix} gives the implementation-family matrix; do not infer a
cross-renderer guarantee from the common C\texttt{++} type alone.

\section{Model and Stock Effects}
\label{sec:api-model-effects}

\begin{small}
\begin{longtable}{>{\raggedright\arraybackslash}p{3cm}>{\raggedright\arraybackslash}p{5.2cm}>{\raggedright\arraybackslash}p{3.3cm}}
\toprule
\textbf{Type group} & \textbf{Public shape} & \textbf{Boundary} \\
\midrule
\endfirsthead
\toprule
\textbf{Type group} & \textbf{Public shape} & \textbf{Boundary} \\
\midrule
\endhead
\cnaclass{Model}, \cnaclass{ModelBone}, \cnaclass{ModelMesh}, \cnaclass{ModelMeshPart} & Bone and
  mesh collections, root/parent hierarchy, transform-copy helpers, part buffers/declaration,
  effects and draw. & Four loading routes have different provenance; \texttt{Model::Draw}
  requires matrix-capable effects and uses serial shared scratch state. \\
\cnaclass{BasicEffect} & Matrices, texture/vertex colour, material colours, fog and three
  directional lights. & Renderer-owned shader/fixed-function mapping decides which fields are
  consumed. \\
\cnaclass{AlphaTestEffect} & Texture, vertex colour, fog, alpha comparison and reference value. &
  A draw can succeed while a fallback loses alpha-test semantics. \\
\cnaclass{DualTextureEffect} & Two texture layers, matrices, vertex colour and fog. & The second
  layer needs a renderer route and compatible packed layout. \\
\cnaclass{EnvironmentMapEffect} & Diffuse texture, cube environment map, amount/specular/fresnel,
  matrices, fog and lights. & Cube construction/sampling and effect dispatch are separate gates. \\
\cnaclass{SkinnedEffect} & Matrices/material/texture/fog/lights plus bone palette and 1/2/4
  weights per vertex. & Maximum palette/layout and animation binding are runtime contracts; CNA's
  vertex-colour addition is marked \texttt{CNAEXT}. \\
\bottomrule
\end{longtable}
\end{small}

Chapter~\ref{ch:models-meshes} covers runtime ownership, Chapters~\ref{ch:gltf-import-core}--
\ref{ch:gltf-conformance} cover glTF/CNJ, and Appendix~\ref{ch:gltf-evidence-matrix} keeps
importer, runtime and test evidence separate.

\section{GraphicsDevice and State Objects}
\label{sec:api-graphicsdevice}

\cnaclass{GraphicsDevice} is the renderer-independent dispatch facade.  Its public groups are:

\begin{itemize}
  \item adapter/display/profile/status/presentation identity and disposal state;
  \item texture, sampler, vertex-buffer, index-buffer and render-target binding;
  \item blend, depth/stencil, rasterizer, viewport, scissor and related scalar state;
  \item clear, present, backbuffer readback and screenshot helpers;
  \item bound-buffer, user-array, indexed and instanced draw families;
  \item capability queries, unsupported-3D policy, and marked diagnostic/recovery helpers.
\end{itemize}

The facade stores and returns state even where an implementation family does not consume it.
The pin also contains concrete asymmetries in singular/plural bindings, offsets/frequencies,
readback source, clear masks, presentation interval/format and capability truth.  Therefore use
the operation-level matrices in Chapters~\ref{ch:graphicsdevice} and~\ref{ch:renderer-contract};
\texttt{SupportsCapability(true)} is a preflight hint, not an execution artifact.

The four state objects remain distinct:

\begin{center}
\begin{tabular}{p{3.2cm}p{8cm}}
\toprule
\textbf{Type} & \textbf{State represented} \\
\midrule
\cnaclass{BlendState} & colour/alpha factors and equations, write masks, blend factor and presets \\
\cnaclass{DepthStencilState} & depth enable/write/compare and front/back stencil operations \\
\cnaclass{RasterizerState} & cull/fill, bias, multisampling and scissor-enable intent \\
\cnaclass{SamplerState} & addressing, filtering, anisotropy, mip limit and LOD bias \\
\bottomrule
\end{tabular}
\end{center}

CNA copies assigned state objects by value.  Later mutation of the original object does not
retroactively modify the device copy.  Whether each copied field reaches native state still
depends on the renderer family.

\section{Effect Base, Parameters, and Queries}
\label{sec:api-buffers-effect-base}

\cnaclass{Effect} owns parameter and technique collections, a current technique, cloning, and
the CNA-added convenience that applies the current pass.  Its bytecode constructor accepts the
XNA/FNA D3D9 Effect Framework payload supported by the selected renderer; FNA3D is the normal
alpha.1 route, while SDL\_GPU, EasyGL and Vulkan require their explicit experimental effect
options.  \cnaclass{EffectParameter} exposes
typed scalar/vector/matrix/string/texture setters and getters plus class/type/dimension metadata
and nested collections.  This is not a source compiler and does not accept \texttt{.fx}/HLSL,
DXBC or MGFX.  \cnaclass{ShaderEffect} is a distinct renderer-specific source/program route.
The existence of a stock effect or parameter setter therefore does not prove compiled-effect
execution for a renderer absent the matching capability and factory.

\cnaclass{OcclusionQuery} exposes \texttt{Begin}, \texttt{End}, completion, and pixel count.
Appendix~\ref{ch:feature-matrix} identifies families with real, conditional, trace-only, null, or
misreported query paths.  A non-null public wrapper and a true capability bit are insufficient
without a usable renderer object and a discriminating sample-count test.

\section{Reading this appendix safely}

For exact declarations, search the selected type under
\path{modules/<owner>/include/Microsoft/Xna/Framework/} at CNA
\CnaEditionTag{} @ \CnaRevision.  Then pair the header with its implementation,
renderer route, tests and host evidence.  This appendix answers ``where is the surface?''; the
book's chapters and retained audit reports answer ``what does this revision actually do?''
